% Expose IX, section 10 opening and sections 10.1--10.2.
% French transcription lines 2717--2786; authority scan physical pp. 451--453
% (scan indices 450--452; source folios 439--441). Physical p. 454 was
% checked for the exact 10.3 boundary. All pages were compared in four
% overlapping 1100-dpi bands; the print is clear at that resolution.

\medskip
\noindent
10. \underline{Properties of the monodromy pairing}. \underline{Integrality
and positivity theorem}

In this section we shall first give certain essentially trivial formal
properties of the monodromy pairings, leaving the (somewhat tedious)
details of the verifications to the reader. We shall also state a deeper
theorem (10.4) concerning these pairings and indicate how its proof may be
reduced to a special case of Jacobians, to be treated (by transcendental
methods) in section 12, which is in fact logically independent of section
11.

\medskip
\noindent
10.1. \underline{Transport of structure}. It is clear from the definition
of the monodromy pairing (9.1.2) that this pairing is compatible with
transport of structure in the situation $(S,A_K,A'_K)$.

\medskip
\noindent
10.2. \underline{Symmetry}. Starting with the dual scheme $A'_K$ of $A_K$
and applying the construction of section 9 to it, we must introduce the
dual scheme $(A'_K)'$ of $A'_K$, which is canonically isomorphic to $A_K$
by the biduality theorem ([5] or VIII 3.2). Denote by
$u_\ell^{A_K}$ the monodromy pairing (9.1.2) relative to $A_K$. In the same
way we obtain a pairing

\begin{equation}\tag{10.2.1}
u_\ell^{A'_K}:\underline M'_\ell\otimes\underline M_\ell
 \longrightarrow\mathbb Z_\ell.
\end{equation}
I claim that the latter is obtained from $u_\ell^{A_K}$ by the symmetry

\begin{equation*}
s:\underline M'_\ell\otimes\underline M_\ell
 \xrightarrow{\ \sim\ }
 \underline M_\ell\otimes\underline M'_\ell;
\end{equation*}
that is,

\begin{equation}\tag{10.2.2}
u_\ell^{A'_K}=u_\ell^{A_K}\circ s.
\end{equation}

Taking account of the known anticommutativity property of the pairings
$\varphi_\xi$ ([5] or VIII 2.2.11), the symmetry formula (10.2.2) may be
reformulated as follows in terms of a fixed polarization $\xi$ of $A_K$.
It defines a homomorphism

\begin{equation*}
\widetilde\xi:A_K\longrightarrow A'_K
\end{equation*}
and hence a corresponding homomorphism ($t$ standing for ``toric'')

\begin{equation}\tag{10.2.3}
\xi_t:\underline M\longrightarrow\underline M'.
\end{equation}
We may therefore form the pairing

\begin{equation}\tag{10.2.4}
\begin{aligned}
\varphi^t_{\xi,\ell}:\underline M_\ell\otimes\underline M_\ell
 &\longrightarrow\mathbb Z_\ell,\\
\varphi^t_{\xi,\ell}(x,y)&=u_\ell^{A_K}\bigl(x,\xi_t(y)\bigr).
\end{aligned}
\end{equation}
The preceding form $\varphi^t_{\xi,\ell}$ is \underline{symmetric}:

\begin{equation}\tag{10.2.5}
\varphi^t_{\xi,\ell}(x,y)=\varphi^t_{\xi,\ell}(y,x).
\end{equation}
In fact, this assertion is valid for every element $\xi$ of the
Néron--Severi group of $A_K$, not necessarily a polarization. But only
when $\xi$ is nondegenerate---that is, defines an \underline{isogeny}
$A_K\longrightarrow A'_K$, and hence also an isogeny (10.2.3)---is the
symmetry relation (10.2.5) equivalent to (10.2.2). When $\xi$ is a
polarization, we shall moreover sharpen (10.2.5) considerably in 10.4(b)
below.

Let us finally mention a third interpretation of the symmetry relation.
Consider a biextension (that is, a divisorial correspondence) $\xi$ on a
pair $(A_K,\overline A_K)$ of abelian schemes with semistable reduction
over $K$. It defines a homomorphism

\begin{equation*}
A_K\longrightarrow\overline A'_K.
\end{equation*}
Passing to the Néron models, this defines a homomorphism on the maximal
tori of the special fibres, and therefore, in the reverse direction, a
homomorphism on their character groups:

\begin{equation}\tag{10.2.6}
\xi_t:\overline{\underline M}\longrightarrow\underline M',
\end{equation}
where $\overline{\underline M}$ is defined from $\overline A_K$ as
$\underline M$ is defined from $A_K$. We can therefore form the pairing

\begin{equation}\tag{10.2.7}
\begin{aligned}
\varphi^t_\xi:\underline M_\ell\otimes
 \overline{\underline M}_\ell&\longrightarrow\mathbb Z_\ell,\\
\varphi^t_\xi(x,y)&=u_\ell^{A_K}\bigl(x,\xi_t(y)\bigr).
\end{aligned}
\end{equation}
\footnote{In the printed right-hand side the underlying pairing is denoted
only by $\varphi^t$, without a definition. It is the monodromy pairing
$u_\ell^{A_K}$ used in the identical construction (10.2.4).}
Notice in passing that, for a variable abelian scheme $A_K$ with semistable
reduction, the twisted constant group $\underline M$ is a covariant
functor of $A_K$. Formation of the pairings (10.2.7) associated with
biextensions of abelian schemes with semistable reduction over $(K,S)$ is
then compatible, in the evident sense, with taking inverse images of
biextensions.

When $\xi$ is the canonical (``Weil'') biextension of $(A_K,A'_K)$, the
pairing (10.2.7) is of course precisely the monodromy pairing (9.1.2).
\footnote{The printed source refers here to (9.2.1), which is the inertia
orthogonality formula $gx-x\in V^\perp$, not the monodromy pairing. The
pairing is numbered (9.1.2).}
Now consider the transpose ${}^s\xi$ of $\xi$, a biextension of
$(\overline A_K,A_K)$, and hence the pairing

\begin{equation*}
\varphi^t_{{}^s\xi}:\overline{\underline M}_\ell\otimes
 \underline M_\ell\longrightarrow\mathbb Z_\ell.
\end{equation*}
The symmetry relation is

\begin{equation}\tag{10.2.8}
\varphi^t_{{}^s\xi}=\varphi^t_\xi\circ s,
\qquad
s:\overline{\underline M}_\ell\otimes\underline M_\ell
 \xrightarrow{\ \sim\ }
 \underline M_\ell\otimes\overline{\underline M}_\ell,
\end{equation}
where $s$ is again the symmetry isomorphism. Compare this relation with
the antisymmetry relation of VIII 2.2.11.

